\begin{lemma}[the cusp boundary defect and Chern-class specialization]
\hypertarget{tag-03ID}{}\label{wbII:lem:picard-cusp-defect}
\StatementRecord{wbII:lem:picard-cusp-defect}{proved}{03ID}{\SeriesRef{wbII:lem:picard-crossed-defect}, \SeriesRef{wbII:eq:W:5.17g}, \SeriesRef{wbII:eq:W:5.17j'}, \SeriesRef{wbII:eq:W:5.17p}}

\ In the notation of
\cref{wbII:lem:picard-crossed-defect}, one has
\[
 \bar d_0=0\quad\text{in }V_G\otimes\mathbb R.
\tag{W.36}\hypertarget{tag-03H3}{}\label{wbII:eq:W:5.17n}
\]
Moreover, for every line bundle $L$ on the cusp neighbourhood $N_0$,
nearby-cycle specialization is compatible with first Chern classes:
\[
 \operatorname{sp}^2\bigl(c_1(L|_W)\bigr)=c_1(L|_F).
\tag{W.37}\hypertarget{tag-03H4}{}\label{wbII:eq:W:5.17o}
\]
\end{lemma}



\begin{proof}

\begin{enumerate}[label={$\langle1\rangle$\arabic*.},leftmargin=*]
\item Let $\widetilde a_0$ be the flat-marked clockwise zero-section meridian in
the lifted relation of \eqref{wbII:eq:W:5.17g}, and let
\[
 \xi_{\widehat\delta}
 =a_*([S^1_{\widehat\delta}]\times[\widetilde a_0])
\]
be its sweep under fibrewise translation by $\widehat\delta$.  The boundary
$\rho_0$ of $P$ was chosen to be its projection.  This is the zero-section
sweep: there is no hidden vertical drift in $\widetilde a_0$, by the
explicit overlap calculation in \CompendiumRef{I}{mapsI:top:prop:geometric-pi1}.  Hence the
vertical form $\Omega_E$ vanishes on the swept torus, and the slant identity
\eqref{wbII:eq:W:5.17j'} gives
\[
 \langle d_0,\widehat\delta\rangle
 =-\langle c_1(L|_{J_P})-[\Omega_E],
          \xi_{\widehat\delta}\rangle
 =-\langle c_1(L|_{J_P}),\xi_{\widehat\delta}\rangle.
\tag{W.38}\hypertarget{tag-03H5}{}\label{wbII:eq:W:5.17p}
\]
Combine the exponential overlap coordinates $(x_1,x_2,t)$ given by $x_k=e^{2\pi i\zeta_k}$ and $t=e^{2\pi is}$ with
the unimodular cusp chart
\[
 (u,x_1,x_2)=\left(\frac{t}{x_1x_2},x_1,x_2\right).
\]
The literal solid torus
\[
 \mathbb S_{\widehat\delta}
 =\{|u|=\rho,\ x_1=1,\ |x_2|\le1\}
\]
has boundary parametrization
\[
 (\theta,\phi)\longmapsto
 \bigl(\rho e^{i(\theta-\phi)},1,e^{i\phi}\bigr)
\tag{W.39}\hypertarget{tag-03H7}{}\label{wbII:eq:W:5.17p0}
\]
and hence $t=\rho e^{i\theta}$; its $\phi$-circle is translation by
$\widehat\delta$.  Direct evaluation in the proved cusp marking therefore
identifies $\partial\mathbb S_{\widehat\delta}$ with
$\xi_{\widehat\delta}$.  Reversing $\theta$ supplies the fixed clockwise
convention and only changes the common orientation sign.  No reduced
Mayer--Vietoris block or Leray comparison is used here.
Since $L$ extends across $N_0$, its Chern class evaluates to zero on that
boundary.  Thus the left side of \eqref{wbII:eq:W:5.17p} is zero.  By
\CompendiumRef{I}{wb:thm:coinvariants}, $V_G=\mathbb Z[\delta]$, and evaluation on the
globally invariant dual vector $\widehat\delta$ identifies
$V_G\otimes\mathbb R$ with $\mathbb R$.  This proves \eqref{wbII:eq:W:5.17n}.

\item For \eqref{wbII:eq:W:5.17o}, let $\ell=c_1(L|_{N_0})$, let
$\iota_0:W\hookrightarrow N_0$, and let
$\iota_t:F_t\hookrightarrow N_0$ be a sufficiently near smooth cusp fibre.
The neighbourhood realization of the nearby-cycle unit
$\mathbb Z_W\to\psi_{f_0}\mathbb Z$ recalled in the proof of
\CompendiumRef{I}{mapsI:top:thm:p0-integral-specialization} gives the commutative square
\[
\begin{CD}
 H^2(N_0;\mathbb Z) @>{\iota_0^*}>> H^2(W;\mathbb Z)\\
 @V{\iota_t^*}VV @VV{\operatorname{sp}^2}V\\
 H^2(F_t;\mathbb Z) @= H^2(F_t;\mathbb Z).
\end{CD}
\tag{W.40}\hypertarget{tag-03H6}{}\label{wbII:eq:W:5.17p'}
\]
Applying this square to $\ell$ gives
$\operatorname{sp}^2(c_1(L|_W))=c_1(L|_{F_t})$.  Flat parallel transport
in the marked smooth family identifies the latter with $c_1(L|_F)=E$.
Equivalently, $c_1$ is the connecting morphism in the exponential sequence,
and restriction to the shrinking neighbourhoods and to $F_t$ commutes with
that morphism before passage to the direct limit.  This proves the claimed
naturality without using any global Leray differential.

\end{enumerate}
\end{proof}